% fig-ch0-bdi-execution.tex -- redrawn to the house figure language (v2).
% Reader question: how does a BDI agent choose what to do each cycle, and
% what happens when an action fails?
% Claim: events enter the event queue E; S_E selects one event; relevant
% plans are filtered to applicable ones against the belief base bs; S_O
% selects a plan, pushed on the intention stack is; S_I selects which
% intention advances and executes one action a in the host environment; a
% failed action posts a failure event -!g, which is added back to E and
% (per the SOS rule) pops the failed frame and drives recovery.
% Grammar chosen: a cycle on a weighted spine, S_E/S_O/S_I as numbered
% badges keyed to a ruled legend, the belief base as a subordinate artifact
% feeding the filter step, success closing the loop on the left rail and
% failure (pd breach, dashed) closing it on the right rail.  Rejected: the
% "paradigm spectrum" panel (reactive / planning / BDI) from the prior
% draft -- that is a 3-way classification, which is a table, not a figure
% (see report); drawing it here would duplicate the cycle's own claim.
% Facts: Bordini et al. 2007 (AgentSpeak(L) SOS); whitepaper/chapter-0-prerequisites.tex
% Sec. "Operational Semantics of AgentSpeak(L)", the S_E/S_O/S_I definitions
% and the -!g failure rule immediately above \ref{fig:bdi-execution-flow}.
\pdfigurehue{pdcobalt}%
\begin{tikzpicture}[pd figure,x=1cm,y=1cm]
  % ---- one grid, declared once --------------------------------------------
  \def\dy{1.45}    % spine pitch
  \def\sx{6.80}    % spine column
  \def\gx{4.95}    % right edge of the naming gutter (badges sit here)
  \def\lx{1.15}    % left rail: the success loop (perceive)
  \def\ax{10.30}   % right column: belief base and recovery plan
  \def\rx{12.60}   % right rail: the failure loop (-!g re-enters E)
  \def\R{12.60}    % right margin of the legend rule
  % ---- the honest cycle, in the chapter hue -------------------------------
  \node[pd focus state,minimum width=3.3cm] (E)  at (\sx,{4*\dy}) {event queue $\mathcal{E}$};
  \node[pd focus state,minimum width=3.3cm] (AP) at (\sx,{3*\dy}) {applicable plans};
  \node[pd focus state,minimum width=3.3cm] (IS) at (\sx,{2*\dy}) {intention stack};
  \node[pd focus state,minimum width=3.3cm] (XA) at (\sx,{1*\dy}) {execute $a$};
  \foreach \a/\b in {E/AP,AP/IS,IS/XA}{\draw[pd focus arrow] (\a)--(\b);}
  % ---- every selection function named once, in one gutter ----------------
  \foreach \i/\t in {3.5/select event, 2.5/select plan, 1.5/select intention}{
    \node[pd label,anchor=east] at ({\gx-.35},{\i*\dy}) {\t};}
  \foreach \i/\n in {3.5/1, 2.5/2, 1.5/3}{\node[pd badge] at (\gx,{\i*\dy}) {\n};}
  % ---- the belief base: a subordinate feed into the filter step ----------
  \node[pd artifact,minimum width=2.6cm] (BS) at (\ax,{3*\dy}) {belief base $bs$};
  \draw[pd thin arrow] (BS.west) -- (AP.east);
  % ---- the one branch that leaves the honest path: action failure --------
  \node[pd breach state,minimum width=2.6cm] (RP) at (\ax,{1*\dy}) {recovery plan};
  \draw[pd breach arrow] (XA.east) -- (RP.west);
  \node[pd label,text=pderror] at ($(XA.east)!.5!(RP.west)+(0,.60)$) {$-!g$};
  \draw[pd breach arrow,rounded corners=3pt]
    (RP.east) -- (\rx,{1*\dy}) -- (\rx,{4*\dy}) -- (E.east);
  \node[pd badge] at (\rx,2.9) {4};
  % ---- success: the action's perceptions close the loop on the left ------
  \draw[pd focus arrow,rounded corners=3pt]
    (XA.west) -- (\lx,{1*\dy}) -- (\lx,{4*\dy}) -- (E.west);
  \node[pd focus label,rotate=90] at ({\lx-.48},{2.5*\dy}) {perceive};
  % ---- the legend: each selection function, what it picks ----------------
  \def\ly{0.30}
  \draw[pd rule] (0,\ly) -- (\R,\ly);
  \node[pd title,anchor=north west] at (.05,{\ly-.30}) {function};
  \node[pd title,anchor=north west] at (3.90,{\ly-.30}) {selects};
  \draw[pd hairline] (0,{\ly-.95}) -- (\R,{\ly-.95});
  \foreach \i/\n/\f/\s in {%
    0/1/{$S_E(\mathcal{E})$}/{one event: external perception or added goal},
    1/2/{$S_O(\mathrm{App})$}/{one plan, filtered against $bs$, pushed on $is$},
    2/3/{$S_I(is)$}/{which intention advances; runs one action this cycle}}{
      \node[pd badge] at (.28,{\ly-1.35-\i*.62}) {\n};
      \node[pd label,anchor=base west] at (.62,{\ly-1.35-\i*.62}) {\f};
      \node[pd label,anchor=base west,align=left] at (3.90,{\ly-1.35-\i*.62}) {\s};}
  \node[pd badge] at (.28,{\ly-1.35-3*.62}) {4};
  \node[pd label,anchor=base west,text=pderror] at (.62,{\ly-1.35-3*.62}) {$-!g$ (failure)};
  \node[pd label,anchor=base west,align=left,text=pderror] at (3.90,{\ly-1.35-3*.62})
    {pops the failed frame; posts recovery back to $\mathcal{E}$};
  \node[pd note,anchor=north west,text width=11.6cm] at (.05,{\ly-1.35-3*.62-.55})
    {both the honest and the breach path re-enter $\mathcal{E}$: a perception as the
    next event, $-!g$ as the next recovery event.};
\end{tikzpicture}
